\pdfminorversion=7%
\documentclass[aspectratio=169,mathserif,notheorems]{beamer}%
%
\xdef\bookbaseDir{../../bookbase}%
\xdef\sharedDir{../../shared}%
\RequirePackage{\bookbaseDir/styles/slides}%
\RequirePackage{\sharedDir/styles/styles}%
\toggleToGerman%
%
\subtitle{26.~yEd Installieren}%
%
\begin{document}%
%
\startPresentation%
%
\section{Einleitung}%
%
\begin{frame}[b]%
\frametitle{Einleitung}%
\begin{itemize}%
\item Ein wichtiger Schritt im Datenbankentwicklungsprozess ist das Erstellen von abstrakten konzeptuellen Modelle der Problemdomäne.%
\item<2-> Dazu werden graphische Diagramme, sogenannte Entity-Relationship Diagramme~(\glsShort{ERD}s) gezeichnet, die die Objekte und deren Beziehungen aus der realen Welt darstellen, die wir in der Datenbank modellieren wollen.%
\item<3-> Diese konzeptuellen Modelle sind unabhängig von der letztendlich verwendeten \glsShort{dbms}-Technologie.%
\item<4-> Als Werkzeug um solche Diagramme zu zeichnen verwenden wir \yEd\cite{SG2015MDAWY,Y2011YGEM}.%
\end{itemize}%
\locateGraphic[%
\parbox{0.6\paperwidth}{\strut\\[10pt]The logo of the \href{https://www.yworks.com/products/yed}{\yEd\ graph editor}. %
The \yEd~logo is protected by copyright. %
\yEd~is a registered trademark of \href{https://www.yworks.com}{yWorks~GmbH}. %
Unauthorized use, reproduction, or distribution is strictly prohibited.}%
]{}{width=0.2\paperwidth}{../05_software_und_literatur/graphics/yEdLogo}{0.6}{0.06}%
\end{frame}%
%
\section{Ubuntu Linux}%
%
\begin{frame}[t]%
\frametitle{Ubuntu Linux}%
\begin{itemize}%
%
\only<-2>{%
\item Wir werden \yEd\ von der Webseite \url{https://www.yworks.com/products/yed} herunterladen.%
\item<2-> Auf dieser Webseite klicken wir auf den \menu{Download}-Button.%
}%
%
\only<-3>{%
\item<3-> Dann klicken wir auf \emph{More yEd downloads}\dots%
}%
%
\only<-4>{%
\item<4-> Und klicken wir dann auf \emph{Show all yEd downloads}.%
}%
%
\only<-5>{%
\item<5-> Wir wollen \emph{Zipped yEd \textil{jar} file} herunterladen, was für alle Systeme funktioniert, auf denen schon \pgls{Java} installiert ist.%
}%
%
\only<-6>{%
\item<6-> Falls \pgls{Java} auf Ihrem Computer noch nicht installiert ist, instalieren Sie es via \bashil{sudo apt-get install openjdk-xx-jre}, wobei \textil{xx} mit der aktuellen Version ersetzt werden kann, \DEzB~\textil{8} oder \textil{21}.%
%
\item<7-> Wir klicken also auf \menu{Download .zip file}.%
}%
%
\only<-9>{%
\item<8-> Nachdem wir auf den Download-Button geklickt habe, kommen wir zum Lizentbildschirm.%
%
\item<9-> Nachdem wir die Lizenz gelesen haben und wenn wir mit ihr OK sind, stimmen wir mit \menu{I accept the license terms} zu.%
}%
%
\only<-10>{%
\item<10-> Nach dem wir die Lizenz angenommen haben, können wir auf \menu{Download} klicken.%
}%
%
\only<-11>{%
\item<11-> Der Download beginnt.%
}%
%
\only<-13>{%
\item<12-> Irgendwann ist der Download fertig.%
%
\item<13-> Wir können auf den Ordner klicken, in den die Datei heruntergeladen wurde.%
}%
%
\only<-15>{%
\item<14-> Wir finden die \textil{zip}-Datei in diesem Ordner.%
\item<15-> Wir entpacken den Order in das Verzeichnis, in das wir \yEd\ installieren/benutzen wollen.%
}%
%
\only<-18>{%
\item<16-> Wir öffnen ein \glslink{terminal}{Terminal}, in dem wir \ubuntuTerminal\ drücken.%
\item<17-> Wir gehen in den Ordner, in dem die ausgepackten Dateien sind und finden die Datei \bashil{yed.jar}.%
\item<18-> In meinem Fall ist das \bashil{/tmp/yed} {\dots} Sie installieren das natürlich woanders hin.%
}%
%
\only<-19>{%
\item<19-> Wir staten das Programm mit dem Kommando \bashil{java -jar yed.jar}.%
}%
%
\only<-20>{%
\item<20-> Der Spash-Screen erscheint, während das Programm lädt.%
}%
%
\only<-21>{%
\item<21-> Das Programm ist gestartet.%
}%
%
\end{itemize}%
%
\locateGraphicTB{1-2}{width=0.45\paperwidth}{graphics/installingYedUbuntu/installYedUbuntu01website}{0.275}{0.33}%
\locateGraphicTB{3}{width=0.5\paperwidth}{graphics/installingYedUbuntu/installYedUbuntu02downloadsA}{0.25}{0.33}%
\locateGraphicTB{4}{width=0.5\paperwidth}{graphics/installingYedUbuntu/installYedUbuntu03downloadsB}{0.25}{0.33}%
\locateGraphicTB{5-7}{width=0.5\paperwidth}{graphics/installingYedUbuntu/installYedUbuntu04downloadsC}{0.25}{0.33}%
\locateGraphicTB{8-9}{width=0.5\paperwidth}{graphics/installingYedUbuntu/installYedUbuntu05license}{0.25}{0.33}%
\locateGraphicTB{10}{width=0.5\paperwidth}{graphics/installingYedUbuntu/installYedUbuntu06licenseOkDownload}{0.25}{0.33}%
\locateGraphicTB{11}{width=0.55\paperwidth}{graphics/installingYedUbuntu/installYedUbuntu07downloading}{0.225}{0.33}%
\locateGraphicTB{12-13}{width=0.55\paperwidth}{graphics/installingYedUbuntu/installYedUbuntu08downloaded}{0.225}{0.33}%
\locateGraphicTB{14-15}{width=0.55\paperwidth}{graphics/installingYedUbuntu/installYedUbuntu09inNautilus}{0.225}{0.33}%
\locateGraphicTB{16-18}{width=0.64\paperwidth}{graphics/installingYedUbuntu/installYedUbuntu10terminalInFolder}{0.18}{0.4}%
\locateGraphicTB{19}{width=0.64\paperwidth}{graphics/installingYedUbuntu/installYedUbuntu11javaJar}{0.18}{0.4}%
\locateGraphicTB{20}{width=0.5\paperwidth}{graphics/installingYedUbuntu/installYedUbuntu12splash}{0.25}{0.33}%
\locateGraphicTB{21}{width=0.64\paperwidth}{graphics/installingYedUbuntu/installYedUbuntu13yEd}{0.18}{0.33}%
\end{frame}%
%
%
\section{Microsoft Windows}%
%
\begin{frame}[t]%
\frametitle{Microsoft Windows}%
\begin{itemize}%
%
\only<-2>{%
\item Wir werden \yEd\ von der Webseite \url{https://www.yworks.com/products/yed} herunterladen.%
\item<2-> Auf dieser Webseite klicken wir auf den \menu{Download}-Button.%
}%
%
\only<-3>{%
\item<3-> Wir finden den \menu{Download}-Button und klicken darauf.%
}%
%
\only<-4>{%
\item<4-> Wir klicken auf \menu{yEd for Windows}.%
}%
%
\only<-6>{%
\item<5-> Nachdem wir auf den Download-Button geklickt habe, kommen wir zum Lizentbildschirm.%
%
\item<6-> Nachdem wir die Lizenz gelesen haben und wenn wir mit ihr OK sind, stimmen wir mit \menu{I accept the license terms} zu.%
}%
%
\only<-7>{%
\item<7-> Nach dem wir die Lizenz angenommen haben, können wir auf \menu{Download} klicken.%
}%
%
\only<-8>{%
\item<8-> Der Download beginnt.%
}%
%
\only<-9>{%
\item<9-> Nachdem der Download fertig ist, klicken wir auf \menu{Open file} oder etwas äquivalentes in Ihrem Browser.%
}%
%
\only<-10>{%
\item<10-> Die Installation beginnt.%
}%
%
\only<-12>{%
\item<11-> \microsoftWindows\ fragt uns, ob wir der Installation erlauben wollen, Änderungen an unserem Computer vorzunehmen.%
\item<12-> Wir klicken~\menu{Yes}.%
}%
%
\only<-13>{%
\item<13-> Wir klicken \menu{Next} im Willkommensbildschirm.
}%
%
\only<-15>{%
\item<14-> Um \yEd\ zu installieren, müssen wir der Lizenz zustimmen.%
\item<15-> Wenn Sie mit der Lizenz einverstanden sind, können Sie ihr zustimmen und \menu{Next} klicken.%
}%
%
\only<-16>{%
\item<16-> Wir ändern den Installationsordner nicht und klicken \menu{Next}.%
}%
%
\only<-17>{%
\item<17-> Wir ändern auch die Startmenüeinstellungen nicht und klicken \menu{Next}.%
}%
%
\only<-18>{%
\item<18-> Wir ändern auch die Zuständigkeiten für Dateitypen nicht und klicken \menu{Next}.%
}%
%
\only<-19>{%
\item<19-> Wir ändern auch die Einstellungen für Icons nicht und klicken \menu{Next}.%
}%
%
\only<-20>{%
\item<20-> Die Installation beginnt.%
}%
%
\only<-22>{%
\item<21-> Die Installation ist fertig.%
\item<22-> Wir markieren \inQuotes{Run yEd Graph Editor} und klicken auf~\menu{Finish}.%
}%
%
\only<-23>{%
\item<23-> Der Splash-Screen erscheint.%
}%
%
\only<-24>{%
\item<24-> \yEd\ läuft!%
}%
\end{itemize}%
%
\locateGraphicTB{1-2}{width=0.64\paperwidth}{graphics/installingYedWindows/installingYedWindows01website}{0.18}{0.33}%
\locateGraphicTB{3}{width=0.64\paperwidth}{graphics/installingYedWindows/installingYedWindows02download}{0.18}{0.33}%
\locateGraphicTB{4}{width=0.64\paperwidth}{graphics/installingYedWindows/installingYedWindows03downloadForWindows}{0.18}{0.33}%
\locateGraphicTB{5-6}{width=0.64\paperwidth}{graphics/installingYedWindows/installingYedWindows04license}{0.18}{0.33}%
\locateGraphicTB{7}{width=0.64\paperwidth}{graphics/installingYedWindows/installingYedWindows05acceptLicenseAndDownload}{0.18}{0.33}%
\locateGraphicTB{8}{width=0.64\paperwidth}{graphics/installingYedWindows/installingYedWindows06downloading}{0.18}{0.33}%
\locateGraphicTB{9}{width=0.64\paperwidth}{graphics/installingYedWindows/installingYedWindows07downloaded}{0.18}{0.33}%
\locateGraphicTB{10}{width=0.64\paperwidth}{graphics/installingYedWindows/installingYedWindows08preparingInstall}{0.18}{0.33}%
\locateGraphicTB{11-12}{width=0.5\paperwidth}{graphics/installingYedWindows/installingYedWindows09shouldWeInstall}{0.25}{0.33}%
\locateGraphicTB{13}{width=0.4\paperwidth}{graphics/installingYedWindows/installingYedWindows10welcome}{0.3}{0.33}%
\locateGraphicTB{14}{width=0.4\paperwidth}{graphics/installingYedWindows/installingYedWindows11licenseAcceptYesNo}{0.3}{0.33}%
\locateGraphicTB{15}{width=0.4\paperwidth}{graphics/installingYedWindows/installingYedWindows12licenseAccept}{0.3}{0.33}%
\locateGraphicTB{16}{width=0.4\paperwidth}{graphics/installingYedWindows/installingYedWindows13destination}{0.3}{0.33}%
\locateGraphicTB{17}{width=0.4\paperwidth}{graphics/installingYedWindows/installingYedWindows14startMenu}{0.3}{0.33}%
\locateGraphicTB{18}{width=0.4\paperwidth}{graphics/installingYedWindows/installingYedWindows15fileSuffixes}{0.3}{0.33}%
\locateGraphicTB{19}{width=0.4\paperwidth}{graphics/installingYedWindows/installingYedWindows17icon}{0.3}{0.33}%
\locateGraphicTB{20}{width=0.4\paperwidth}{graphics/installingYedWindows/installingYedWindows18installing}{0.3}{0.33}%
\locateGraphicTB{21-22}{width=0.4\paperwidth}{graphics/installingYedWindows/installingYedWindows19doneRun}{0.3}{0.33}%
\locateGraphicTB{23}{width=0.55\paperwidth}{graphics/installingYedWindows/installingYedWindows20splash}{0.225}{0.33}%
\locateGraphicTB{24}{width=0.5\paperwidth}{graphics/installingYedWindows/installingYedWindows21running}{0.25}{0.33}%
\end{frame}%
%
\section{Zusammenfassung}%
%
\begin{frame}%
\frametitle{Zusammenfassung}%
\begin{itemize}%
\item Wir haben nun \yEd\ installiert.%
\item<2-> In sehr vielen Datenbankvorlesungen werden \glslink{ERD}{ERDs} eingeführt und ausführlich diskutiert.%
\item<3-> Die Studenten müssen solche dann in Prüfungen malen können.%
\item<4-> Aber wie malt man ERDs im richtigen Leben?%
\item<5-> Sicher nicht auf Papier.%
\item<6-> Dann versuchen die Studenten später, ERDs mit Grafikprogrammen wie \glsFull{inkscape} zu malen {\dots} und finden sie unpraktisch.%
\item<7-> Weil sie nie die richtigen Werkzeuge kennengelernt haben.%
\item<8-> \yEd\ ist ein richtiges Werzeug.%
\item<9-> Und wenn wir nachher (nach ein paar anderen Diskussionen) anfangen, ERDs zu malen -- dann haben Sie das richtige Werkzeug zur Hand.%
\end{itemize}%
\end{frame}%
%
\endPresentation%
\end{document}%%
\endinput%
%
